\documentclass[10pt,a4paper]{amsart}
\usepackage[margin=19mm]{geometry}
\usepackage{amsmath,amssymb,mathtools}
\usepackage[scr=rsfs,cal=euler]{mathalfa}
\usepackage{xfrac}
\usepackage[unicode,hidelinks,bookmarks=false]{hyperref}
\newcommand{\ie}{i.e.\ }
\newcommand{\cf}{cf.\ }
\newcommand{\resp}{respectively\ }
\newcommand{\Subsection}{\subsection}
\providecommand{\nobreakdash}{\nobreak\hbox{-}}
\setlength{\emergencystretch}{2em}
\multlinegap0pt

\theoremstyle{plain}
\newtheorem{theorem}{Theorem}[section]
\newtheorem{lemma}[theorem]{Lemma}
\newtheorem{proposition}[theorem]{Proposition}
\newtheorem{corollary}[theorem]{Corollary}
\newtheorem{conjecture}[theorem]{Conjecture}
\theoremstyle{remark}
\newtheorem{notation}[theorem]{Notation}
\newtheorem{remark}[theorem]{Remark}
\newtheorem{assumption}[theorem]{Assumption}
\newtheorem{definition}[theorem]{Definition\rm}
\newtheorem{example}{\it Example\/}
\title[Spectral Asymptotics of a Singular Harmonic Oscillator with Aharonov--Bohm Flux]{Spectral Asymptotics of a Singular Harmonic Oscillator with Aharonov--Bohm Flux}
\author{Bernard Helffer, Ayman Kachmar, Fran\c{c}ois Nicoleau}
\date{Source: arXiv:2609.28076v1 (CC BY 4.0)}
\begin{document}
\maketitle

\noindent\emph{Source and licence.} The delivered work is the article shown in the title above, version arXiv:2609.28076v1 (CC BY 4.0), distributed under the Creative Commons Attribution 4.0 International licence, as recorded in the arXiv licence metadata for this exact version and shown on the abstract page saved with this item. The full licence notice, which carries the licence URL and the address of that abstract page, is the bundled file \texttt{licenses/arxiv-2609.28076-CC-BY-4.0.txt}.
\par\smallskip

\noindent\textit{Editorial scope.} Counted unit: the article's Proposition (source label lem:rough-transition) (source0.tex lines 3423--3434, source label \texttt{lem:rough-transition}): ``Let and . Assume that 4kh 1 th , t t 0,t 1 . Then, uniformly for , 0 k(h)-e k(h) O (h 3 2 ).''. It is delivered together with the article's complete original proof (source0.tex lines 3436--3599, one proof environment adjacent to the statement, 164 lines), copied line for line from the article's own text. Required context retained uncounted: eq:1.3 (lines 261--269). Editorial formatting only: an AMS \texttt{amsart} preamble that restates the article's own macro definitions and its own theorem declarations, and the theorem environments the retained lines use; the article's own class file is neither shipped nor input; the retained environments are renumbered sequentially in order of appearance, whereas the article prints them with its own numbering; the article's equation numbering is not reproduced and every cross-reference inside the retained text resolves to the wrapper's numbering. No citation occurs inside any retained line, so the wrapper carries no bibliography; All retained bodies are exact copies of the bundled \texttt{sources/211551/source0.tex}, so no omission receipt applies. No asset-inclusion command occurs inside any retained span and the package ships no image asset.


\section*{Retained context: eq:1.3 (source0.tex lines 261--269)}

\begin{subequations}\label{eq:1.3}
\begin{equation}
\lambda_k(h)\geq e_k(h)\,,
\end{equation}
where 
\begin{equation}
e_k(h):=(4k+2\nu-2)h
\end{equation}
\end{subequations}


\section{The counted unit: Proposition (source label lem:rough-transition) and its complete original proof}

\begin{proposition}\label{lem:rough-transition}
Let $0<\delta<2/3$ and $0<t_0<t_1<+\infty$. Assume that
\[
\rho=4kh=1+th^\delta,
\qquad t\in[t_0,t_1].
\]
Then, uniformly for $t\in[t_0,t_1]$,
\[
0\leq \lambda_k(h)-e_k(h)
=\mathcal O\bigl(h^{3\delta/2}\bigr).
\]
\end{proposition}

\begin{proof}
	%
	By \eqref{eq:1.3}, we have
	\begin{equation}\label{eq:rough-lower-bound}
		e_k(h)\leq\lambda_k(h)\,.
	\end{equation}
	We introduce a Dirichlet decoupling at $x=1$ for the half-line
	operator $T_h^\infty$. The resulting operator is
	\[
	T_h^{\mathrm{dec}}
	=
	T_{h,(0,1)}^{D}\oplus T_{h,(1,\infty)}^{D}\,.
	\]
	The first component is precisely the operator on $(0,1)$ whose
	eigenvalues are $\lambda_j(h)\,$. We denote by
	\[
	\lambda_1^{\mathrm{ext}}(h)
	\leq
	\lambda_2^{\mathrm{ext}}(h)
	\leq\cdots
	\]
	the eigenvalues of the second component, which we shall refer to as
	the exterior component. Thus
	\[
	\sigma\bigl(T_h^{\mathrm{dec}}\bigr)
	=
	\{\lambda_j(h)\}_{j\geq1}
	\cup
	\{\lambda_j^{\mathrm{ext}}(h)\}_{j\geq1}\,,
	\]
	where the union is understood with multiplicities. Let
	\[
	\mu_1(h)\leq\mu_2(h)\leq\cdots
	\]
	denote this union rearranged in increasing order. Since the
	decoupling amounts to imposing one additional Dirichlet condition
	on the form domain, the min--max principle gives
	\begin{equation}\label{eq:rough-interlacing}
		e_j(h)\leq\mu_j(h)\leq e_{j+1}(h)\,.
	\end{equation}
	We next estimate the spectrum of the exterior component. On
	$(1,\infty)\,$, we have
	\[
	x^2+h^2\frac{\nu^2-\frac14}{x^2}
	\geq
	1+2(x-1)-C_\nu h^2\,.
	\]
	Hence, after the scaling $x=1+h^{2/3}y$, the exterior operator is
	bounded from below by
	\[
	1+h^{2/3}A-C_\nu h^2,
	\qquad
	A=-\frac{d^2}{dy^2}+2y\,,
	\]
	on $\mathbb R_+\,$, with Dirichlet boundary condition at $y=0$.
	Let
	\[
	0<\alpha_1<\alpha_2<\cdots
	\]
	denote the eigenvalues of $A$. By the min--max principle,
	\begin{equation}\label{eq:rough-exterior-eigenvalue-lower-bound}
		\lambda_j^{\mathrm{ext}}(h)
		\geq
		1+h^{2/3}\alpha_j-C_\nu h^2\,.
	\end{equation}
	The Airy asymptotics give
	\begin{equation}\label{eq:rough-Airy-growth}
		\alpha_j\sim(3\pi j)^{2/3}\,,
		\qquad j\to+\infty.
	\end{equation}
	By assumption,
	\begin{equation}\label{eq:rough-supercritical-scaling-proof}
		\rho=4kh=1+th^\delta\,.
	\end{equation}
	Set
	\[
	E_h:=1+(t_1+1)h^\delta\,,
	\]
	and let
	\[
	L(h)
	:=
	\#\bigl\{j\geq1:
	\lambda_j^{\mathrm{ext}}(h)\leq E_h\bigr\}\,.
	\]
	By \eqref{eq:rough-exterior-eigenvalue-lower-bound}, if
	$\lambda_j^{\mathrm{ext}}(h)\leq E_h$, then
	\[
	1+h^{2/3}\alpha_j-C_\nu h^2
	\leq
	1+(t_1+1)h^\delta\,.
	\]
	Hence
	\[
	\alpha_j
	\leq
	(t_1+1)h^{\delta-2/3}
	+C_\nu h^{4/3}\,.
	\]
	Using \eqref{eq:rough-Airy-growth}, we infer that
	\begin{equation}\label{eq:rough-L}
		L(h)
		=
		\mathcal O\bigl(h^{3\delta/2-1}\bigr)\,.
	\end{equation}
	We next show that the first $k+L(h)$ eigenvalues of the
	decoupled operator lie below $E_h\,$. Indeed,
	\[
	e_{k+L(h)+1}(h)
	=
	4kh+
	\bigl(4(L(h)+1)+2\nu-2\bigr)h\,.
	\]
	Using \eqref{eq:rough-supercritical-scaling-proof} and
	\eqref{eq:rough-L}, we obtain
	\[
	e_{k+L(h)+1}(h)
	=
	1+th^\delta
	+\mathcal O\bigl(h^{3\delta/2}\bigr)
	+\mathcal O(h)\,.
	\]
	Since $0<\delta<2/3\,$, both remainder terms are $o(h^\delta)$.
	Therefore, uniformly for $t\in[t_0,t_1]\,$,
	\[
	e_{k+L(h)+1}(h)
	\leq
	1+(t_1+1)h^\delta
	=E_h\,,
	\]
	for $h$ sufficiently small. By \eqref{eq:rough-interlacing},
	\[
	\mu_{k+L(h)}(h)
	\leq
	e_{k+L(h)+1}(h)
	\leq E_h\,.
	\]
	Thus the first $k+L(h)$ eigenvalues of the decoupled operator lie
	below $E_h$. By the definition of $L(h)\,$, at most $L(h)$ of them
	belong to the exterior spectrum. Hence at least $k$ of them belong
	to the interior spectrum, and therefore
	\[
	\lambda_k(h)
	\leq
	\mu_{k+L(h)}(h)
	\leq
	e_{k+L(h)+1}(h)\,.
	\]
	Together with \eqref{eq:rough-lower-bound}, this gives
	\[
	0\leq\lambda_k(h)-e_k(h)
	\leq
	e_{k+L(h)+1}(h)-e_k(h)
	=
	4(L(h)+1)h\,.
	\]
	Finally, using \eqref{eq:rough-L}, we obtain
	\[
	\lambda_k(h)-e_k(h)
	=
	\mathcal O\bigl(h^{3\delta/2}\bigr)\,,
	\]
	uniformly for $t\in[t_0,t_1]\,$.
\end{proof}

\end{document}
